\section{How to use these files}
\label{sec:org2dc5823}
Each file in this directory covers one feature area of org.nvim in depth.
They are meant to be \textbf{worked through}, not just read: every section explains
a feature, shows several examples, and ends with exercises.

\begin{itemize}
\item Lines starting with \textbf{Try:} are exercises, with the exact keys to press.
\item Lines starting with \textbf{Expect:} say what you should see afterwards, so you can
check that the feature works.
\item Lines starting with \texttt{\#} are Org comments that annotate the example next to
them. They are dimmed, and they are never exported.
\item Keys use \texttt{<prefix>} for the org.nvim prefix, \texttt{<leader>o} by default.
Press \texttt{g?} in any org buffer to list the keys of that buffer.
\item Nothing breaks if you make a mess: \texttt{u} undoes, and
\texttt{git checkout examples/} restores every file.
\end{itemize}

Start Neovim from the repository root with the bundled init file, so the
examples use a known configuration and your own notes are never touched:

\begin{verbatim}
nvim -u examples/minimal_init.lua examples/00-index.org
\end{verbatim}

It points \texttt{agenda\_files} at the files of this directory and sends captures to
a scratch directory under \texttt{stdpath("state")}. Put the cursor on a link below
and press \texttt{<CR>} to open that file. \texttt{<C-o>} brings you back here.

The examples use dates around late September 2026. When they are in the
past, press \texttt{<C-a>} on the day of a timestamp to move it forward.

If you would rather take a quick tour of everything in one file first, start
with \href{tutorial.org}{tutorial.org}.
\section{The files}
\label{sec:org1183b91}
\subsection{Writing and structure}
\label{sec:orgaf6c356}
\begin{itemize}
\item \href{01-outline.org}{01 Outline}: headlines, folding, motions, moving and sorting subtrees,
narrowing.
\item \href{02-markup.org}{02 Markup}: emphasis, blocks, comments, entities, sub- and
superscripts.
\item \href{03-lists.org}{03 Lists}: plain lists, bullets, counters, checkboxes and progress
cookies.
\end{itemize}
\subsection{Tasks}
\label{sec:org2fc59bc}
\begin{itemize}
\item \href{04-todo.org}{04 TODO items}: keywords, logging, repeaters, habits, dependencies,
priorities.
\item \href{05-tags.org}{05 Tags}: fast selection, groups, inheritance, file tags.
\item \href{06-properties-columns.org}{06 Properties and column view}: drawers, allowed values, inheritance,
column summaries.
\item \href{07-dates.org}{07 Dates and times}: timestamps, SCHEDULED and DEADLINE, the calendar,
date input.
\item \href{08-clocking.org}{08 Clocking}: clocking in and out, effort estimates, clock tables.
\item \href{20-timers-reminders.org}{20 Timers and reminders}: relative and countdown timers, appointment
notifications.
\end{itemize}
\subsection{Finding things}
\label{sec:orge8718bb}
\begin{itemize}
\item \href{09-agenda.org}{09 Agenda}: day and week views, the TODO list, matches, custom commands,
bulk actions.
\item \href{10-sparse-trees.org}{10 Sparse trees}: folding a file down to matches, and the match syntax.
\end{itemize}
\subsection{Collecting and organising}
\label{sec:orgb4ec9e4}
\begin{itemize}
\item \href{11-capture.org}{11 Capture}: templates, expansions, targets.
\item \href{12-refile-archive.org}{12 Refile and archive}: moving subtrees around and out of the way.
\item \href{13-links.org}{13 Links}: every link type, IDs, radio targets, attachments.
\item \href{14-footnotes.org}{14 Footnotes}: inserting, jumping, renumbering.
\end{itemize}
\subsection{Tables and code}
\label{sec:orga368926}
\begin{itemize}
\item \href{15-tables.org}{15 Tables}: creating, aligning, editing, moving, sorting, importing.
\item \href{16-spreadsheet.org}{16 Spreadsheet}: \texttt{\#+TBLFM} formulas, references, functions, formats,
worked sheets.
\item \href{17-babel.org}{17 Babel}: running source blocks, results, variables, noweb, tangling.
\item \href{18-dynamic-blocks.org}{18 Dynamic blocks}: clock tables, column view blocks, custom blocks.
\end{itemize}
\subsection{Publishing and display}
\label{sec:orge87d85b}
\begin{itemize}
\item \href{19-export.org}{19 Export}: the dispatcher, back-ends, export settings.
\item \href{21-images-latex.org}{21 Images and \LaTeX{}}: inline image and \LaTeX{} previews.
\end{itemize}
\subsection{Everything else}
\label{sec:org6852f3b}
\begin{itemize}
\item \href{22-extras.org}{22 Extras}: completion, org-lint, speed keys, inline tasks, encryption,
commands and health checks.
\end{itemize}
